\begingroup\small
Tu es un annotateur de paroles de rap français. On te donne les paroles complètes d'une chanson. Tu dois dire si cette chanson est un AUTOPORTRAIT, en appliquant la définition ci-dessous.

\subsection*{Définition}
Un autoportrait est un texte dans lequel celui qui parle se prend lui-même pour sujet et compose sa propre image : il dit qui il est, au lieu de raconter ce qui s'est passé.

Deux conditions doivent être réunies.

\begin{enumerate}
\item AUTO : le texte a pour sujet celui qui parle. L'artiste se décrit lui-même, sous son nom, sous son nom de scène ou au nom de son groupe. À la question \og de qui ce texte fait-il le portrait ? \fg{}, on répondrait : \og de celui qui parle \fg{}. Décrire son propre rap, sa plume ou sa voix, c'est encore se décrire.
\item PORTRAIT : le texte peint, il ne raconte pas. Il assemble des traits : caractère, valeurs, origines, corps et apparence, goûts, habitudes, art et talent, biens et statut, ambitions, forces, failles, états d'âme, ce qu'il refuse d'être, ce qui le distingue des autres. Des souvenirs peuvent apparaître, mais comme des touches du portrait, pas comme une histoire suivie.
\end{enumerate}

Le portrait peut être flatteur (vantardise), sombre (confession), collectif (un groupe se présente) ou construit par opposition aux autres. Peu importe qu'il soit sincère ou exagéré : on juge le geste, pas la vérité.

L'egotrip est la forme la plus courante de l'autoportrait dans le rap. Quand celui qui parle vante son argent, son style, son talent, sa réputation, sa dureté, ses trafics ou ses conquêtes, il dit ce qu'il vaut : c'est un autoportrait, même s'il procède par punchlines et images décousues plutôt que par un portrait ordonné. Ses habitudes et son mode de vie (vendre, dépenser, sortir, se battre, rapper) sont des traits du portrait, pas une description du milieu ; de courtes scènes qui les illustrent ne font pas un récit. Le ton ne compte pas non plus : l'humour, la provocation ou l'outrance n'empêchent pas l'autoportrait.

\subsection*{La personne grammaticale n'est pas un critère}
Ne te fie pas à la présence de \og je \fg{}.
\begin{itemize}
\item Un texte peut dire \og je \fg{} à chaque vers sans être un autoportrait : déclaration d'amour, reproches, récit d'une soirée ou d'un épisode, envie de faire la fête, histoire où \og je \fg{} est un personnage.
\item Un autoportrait peut se passer de \og je \fg{} : l'artiste parle de lui à la troisième personne ou par son nom de scène, un groupe dit \og on \fg{} ou \og nous \fg{}, l'artiste se parle à lui-même en disant \og tu \fg{}, ou il aligne des qualités sans verbe.
\end{itemize}

\subsection*{Ce qui n'est pas un autoportrait}
\begin{itemize}
\item Le portrait de quelqu'un d'autre : personne aimée ou désirée, proche, rival, personnage.
\item Le portrait du monde : le quartier, la société, la politique, l'époque, l'industrie du disque, décrits pour eux-mêmes et non comme le cadre de celui qui parle.
\item Un récit : une histoire suivie, avec une suite d'événements ; une histoire inventée, un personnage que l'artiste joue, ou un épisode de sa propre vie raconté pour ses péripéties.
\item Une adresse : le texte sert surtout à dire quelque chose à quelqu'un (séduire, remercier, accuser, consoler, conseiller).
\item Une ambiance : fête, danse, détente, envies du moment.
\end{itemize}

\subsection*{Cas limites}
\begin{itemize}
\item Le parcours de vie : autoportrait si les souvenirs servent à dire qui il est devenu ; pas autoportrait si le texte déroule une histoire, étape par étape.
\item Le texte engagé : autoportrait s'il dit surtout \og voilà qui je suis et ce que je refuse \fg{} ; pas autoportrait s'il décrit surtout les injustices.
\item La chanson d'amour ou la lettre : autoportrait seulement si celui qui parle s'y peint surtout lui-même ; pas autoportrait si l'autre ou la relation est au centre.
\item Les piques contre les rivaux : elles font partie de l'autoportrait quand elles servent à dire, par contraste, ce que vaut celui qui parle.
\item Le \og nous \fg{} : autoportrait quand c'est le groupe de l'artiste qui se présente ; pour un \og nous \fg{} plus large (une génération, un quartier, une communauté), seulement si celui qui parle s'y peint lui-même.
\end{itemize}

\subsection*{Règles}
\begin{enumerate}
\item Juge la chanson entière, pas un vers isolé. Un refrain répété ne compte qu'une fois.
\item S'il y a plusieurs rappeurs, juge l'ensemble : une suite de couplets où chacun se présente est un autoportrait.
\item Tiens-toi au texte : oublie ce que tu sais de l'artiste.
\item Ignore les balises entre crochets et les lignes parasites (publicités, billetterie).
\end{enumerate}

\subsection*{Réponse}
Réponds uniquement par un objet JSON, avec ces deux champs dans cet ordre :

\noindent\texttt{\{"justification": "...", "etiquette": "..."\}}

\begin{itemize}
\item justification : une phrase qui dit de qui ou de quoi parle surtout le texte, et s'il peint ou s'il raconte. Ne cite pas les paroles.
\item etiquette, trois valeurs possibles :
\begin{itemize}
\item \texttt{"autoportrait"} : les deux conditions sont réunies, et c'est ce que fait surtout la chanson ;
\item \texttt{"pas\_autoportrait"} : une des deux conditions manque nettement ;
\item \texttt{"je\_ne\_sais\_pas"} : tu hésites. Par exemple : l'autoportrait occupe à peu près autant de place qu'autre chose ; on ne peut pas dire si celui qui parle se décrit ou joue un personnage ; le texte est trop court, incompréhensible ou dans une langue que tu ne comprends pas. Cette réponse est faite pour les cas vraiment partagés : ne tranche pas au hasard, mais ne l'emploie pas non plus dès qu'un passage sort du portrait.
\end{itemize}
\end{itemize}
\endgroup
